% !TEX root = thesis.tex

%% Document Metadata

%% The location of your Assignment Statement (scanned as image)
\thesisspec{figures/zadavaci-list.png}

\title{Final Thesis Template}{Vývoj rozšírenia prehliadača na detekciu toxického obsahu pomocou NLP modelu}
%\thesis{Master thesis}{Diplomová práca}  % typ prace

\author[]{Vladyslav}{Stakhov}[]
\supervisor{Ing. Oliver Lohaj, PhD.}   % thesis supervisor
% \consultant{Donald E. Knuth}  % thesis consultant
% \consultantsaffiliation{Stanford University}{Stanfordská univerzita}  % affiliation of the consultant
% \college{University of Žilina}{Žilinská univerzita}  % univerzita
% \faculty{Faculty of Electrical Engineering and informatics}{Fakulta elektrotechniky a informatiky}  % fakulta
\department{Institute of Artificial Intelligence}{Ústav umeleckej inteligencie}  % katedra
\departmentacr{IAI}{ÚUI}  % skratka katedry
\submissiondate{11}{2}{2026}
\fieldofstudy{Informatika}
\studyprogramme{Informatika}  % iba názov bez číselného kódu
%\city{Košice}  % mesto

\keywords{%
  % english
  browser extension, Chrome Manifest V3, multi-label classification, natural language processing, privacy, TensorFlow.js, toxicity detection
}{%
  % slovak
  Chrome Manifest V3, detekcia toxicity, rozšírenie prehliadača, spracovanie prirodzeného jazyka, súkromie, TensorFlow.js, viacštítková klasifikácia
}

\declaration{
 Vyhlasujem, že som záverečnú prácu vypracoval samostatne s použitím uvedenej odbornej literatúry, ako aj s využitím nástrojov generatívnej umelej inteligencie v súlade s Metodickým pokynom o záverečných a kvalifikačných prácach a Etickým kódexom študenta Technickej univerzity v Košiciach. Umelá inteligencia bola použitá ako podporný nástroj bez porušenia zásad akademickej etiky a autorskej integrity.
}

\abstract{%
  % english
  This thesis focuses on the detection of toxic text in a web browser environment and on the development of a browser extension named \textit{Toxic Text Detector}. The main point of this work is to provide users with a practical tool for identifying potentially harmful online content while preserving user control and their privacy. The solution was designed and implemented as a Chrome Manifest V3 extension. It uses a React-based user interface, a content script for working with web page text, a service worker for communication between components, persistent settings and two inference modes. The local mode does classification directly in the browser, while the remote mode uses an external API through a backend service. Implemented extension can analyze selected text or text elements on a page, mark potentially toxic text, blur it visually and display the detected category together with a confidence score. The solution was evaluated through user interface testing and model experiments focused on local and remote inference. The results of this work show that a hybrid browser-based approach can provide a practical compromise between detection quality, latency, usability and privacy protection. At the same time, the work confirms that toxicity detection remains challenging, especially for minority categories and more context-dependent forms of harmful language.
}{%
  % slovak
  Táto práca sa zameriava na detekciu toxického textu v prostredí webového prehliadača a na vývoj rozšírenia prehliadača s názvom \textit{Toxic Text Detector}. Hlavným cieľom tejto práce je poskytnúť používateľom praktický nástroj na identifikáciu potenciálne škodlivého online obsahu pri zachovaní kontroly používateľov a ich súkromia. Riešenie bolo navrhnuté a implementované ako rozšírenie pre Chrome Manifest V3. Využíva používateľské rozhranie založené na React, skript pre prácu s textom webových stránok, service worker pre komunikáciu medzi komponentmi, trvalé nastavenia a dva režimy inferencie. Lokálny režim vykonáva klasifikáciu priamo v prehliadači, zatiaľ čo vzdialený režim využíva externé API prostredníctvom backendovej služby. Implementované rozšírenie dokáže analyzovať vybraný text alebo textové prvky na stránke, označiť potenciálne toxický text, vizuálne ho rozostriť a zobraziť detegovanú kategóriu spolu s hodnotou spoľahlivosti. Riešenie bolo vyhodnotené prostredníctvom testovania používateľského rozhrania a experimentov s modelmi zameranými na lokálnu a vzdialenú inferenciu. Výsledky tejto práce ukazujú, že hybridný prístup založený na prehliadači môže poskytnúť praktický kompromis medzi kvalitou detekcie, latenciou, použiteľnosťou a ochranou súkromia. Zároveň táto práca potvrdzuje, že detekcia toxicity zostáva náročná, najmä v prípade menšinových kategórií a foriem škodlivého jazyka, ktoré sú viac závislé od kontextu.
}

\acknowledgment{%
  Na tomto mieste by som sa rád poďakoval svojmu vedúcemu práce za jeho čas a odborné vedenie počas riešenia mojej záverečnej práce.
  Rovnako by som sa rád poďakoval svojim rodičom a priateľom za ich podporu a povzbudzovanie počas celého môjho štúdia.
}


% Uncomment the following line for the printer-friendly version of the thesis.
% When defined, final PDF will not contain colored links, which is useful for printing.
% WARNING: Uncomment ONLY when used manually without the Docker image for building the thesis. By uncommenting this line, you will manually override the printable flag and behavior of the Docker image will not work as expected.
% If you want to use the Docker image for building the document, please do not uncomment this line.
%\def\printable{}

